\BourbakiExercises{Exercises}

\begin{enumerate}
\def\labelenumi{\arabic{enumi})}
\item
  Let A be a ring and m a maximal left ideal of A. Let B be the set of elements b of A such that \(mb \subset m\) . Prove that B is the largest subring of A containing m in which m is a two-sided ideal. Let S be the A-module \(A_{s}/m\) , D its commutant, and \(\pi\) the canonical mapping from \(A_{s}\) to S. Prove that, given \(b \in B\) , there exists a unique element \(u_{b}\) of \(\operatorname{End}_{\mathrm{A}}(\mathrm{S})\) such that \(u_{b}(\pi(a)) = \pi(ab)\) for every \(a \in A\) . The mapping \(b \mapsto u_{b}\) induces an isomorphism from B/m to \(\operatorname{End}_{\mathrm{A}}(\mathrm{S})\) .
\item
  Let D be a field, V a right D-vector space, and A a subring of \(\operatorname{End}_{\mathrm{D}}(\mathrm{V})\) that contains the endomorphisms of finite rank (for example, \(A = \operatorname{End}_{\mathrm{D}}(\mathrm{V})\) ). Prove that the dual \(V^{*} = \operatorname{Hom}_{\mathrm{D}}(\mathrm{V}, \mathrm{D})\) is simple as a right A-module.
\item
  Let D be a field, V a right D-vector space, and A the endomorphism ring of V. For every line L in V, we denote by \(m_{L}\) the set of elements u of A such that \(u(L)=0\) . Prove that \(m_{L}\) is a maximal ideal of A, that we have \(m_{L}\neq m_{L'}\) for \(L\neq L'\) , but that the simple A-modules \(A/m_{L}\) are all isomorphic to the A-module V.
\item
  Prove that the Z-module Q admits neither a simple submodule nor a simple quotient module.
\item
  Let \(\mathrm{A}\) be a ring.
\end{enumerate}

\begin{enumerate}
\def\labelenumi{\alph{enumi})}
\tightlist
\item
  Let M, N be monogenous A-modules. Prove that the following properties are equivalent:
\end{enumerate}

\begin{enumerate}
\def\labelenumi{(\roman{enumi})}
\item
  The A-modules M and N are isomorphic.
\item
  For every generator \(m\) of \(\mathrm{M}\) , there exists a generator \(n\) of \(\mathrm{N}\) that has the same annihilator as \(m\) .
\item
  There exist a generator of M and a generator of N that have the same annihilator.
\end{enumerate}

\begin{enumerate}
\def\labelenumi{\alph{enumi})}
\setcounter{enumi}{1}
\item
  Let \(\mathfrak{a},\mathfrak{b}\) be (left) ideals of A. The A-modules A/ \(\mathfrak{a}\) and A/ \(\mathfrak{b}\) are isomorphic if and only if there exists an element \(a\) of A such that \(Aa + \mathfrak{b} = A\) and \(\mathfrak{a}\) is the set of elements \(x\) of A such that \(xa\in \mathfrak{b}\) (use \(a\) ).
\item
  Let \(\mathfrak{m}\) and \(\mathfrak{n}\) be maximal left ideals of A. The simple A-modules A/ \(\mathfrak{m}\) and A/ \(\mathfrak{n}\) are isomorphic if and only if there exists an element \(a\) of A- \(\mathfrak{n}\) such that \(\mathfrak{ma} \subset \mathfrak{n}\) . d) Deduce from b) that the set of quotient A-modules of the A-module \(\mathrm{A}_s\) isomorphic to a given module has cardinal less than \(\operatorname{Card}(\mathrm{A})\) .
\end{enumerate}

\begin{enumerate}
\def\labelenumi{\arabic{enumi})}
\setcounter{enumi}{5}
\item
  Let M be an A-module such that for every \(x \neq 0\) in M, the module Ax is simple. Prove that either M is simple, or the ring of homotheties \(A_{M}\) is a field (if M is not simple, consider two nonzero elements x, y of M such that \(y \notin Ax\) and the annihilator of \(x + y\) ).
\item
  Let S be a simple A-module whose dual \(S^{*}\) and bidual \(S^{**}\) are simple. Prove that the mapping \(u \mapsto {}^{t}u\) is an isomorphism from \(\operatorname{End}_{\mathrm{A}}(\mathrm{S})\) to the opposite field of \(\operatorname{End}_{\mathrm{A}}(\mathrm{S}^{*})\) .
\end{enumerate}

¶ 8) Let A be a ring such that the A-modules \(A_{s}\) and \(A_{d}\) are of finite length. We identify \(A_{d}\) with the dual of \(A_{s}\) . For every left ideal s of A, the orthogonal \(s^{0}\) of s is the right annihilator of s in A; likewise, for every right ideal d, the orthogonal \(d^{0}\) is the left annihilator of d.~Prove that the following properties are equivalent:

\begin{enumerate}
\def\labelenumi{(\roman{enumi})}
\item
  Every simple (left or right) A-module is reflexive.
\item
  The dual of any (left or right) simple A-module is simple.
\item
  For every minimal left ideal \(\mathfrak{s}\) and every minimal right ideal \(\mathfrak{d}\) , we have \(\mathfrak{s}^{00} = \mathfrak{s}\) and \(\mathfrak{d}^{00} = \mathfrak{d}\) .
\item
  The mapping \(\mathfrak{s} \mapsto \mathfrak{s}^0\) is a bijection from the set of left ideals of \(A\) to the set of right ideals.
\end{enumerate}

When these properties hold, we say that the ring A is autoinjective (cf.~X, §1, p.~171, exercice 26).

\begin{enumerate}
\def\labelenumi{\alph{enumi})}
\tightlist
\item
  Suppose that A is autoinjective; let M be a finitely generated (left or right) A-module. Prove that the following properties hold:
\end{enumerate}

(i') The A-module M is reflexive.

(ii') The length of \(\mathbf{M}^*\) is equal to that of \(\mathbf{M}\) .

(iii') For every submodule \(\mathrm{N}\) of \(\mathrm{M}\) , we have \(\mathrm{N}^{00} = \mathrm{N}\) .

(iv') The mapping \(N \mapsto N^{0}\) is a bijection from the set of submodules of M to the set of submodules of \(M^{*}\) .

\begin{enumerate}
\def\labelenumi{\arabic{enumi})}
\setcounter{enumi}{8}
\tightlist
\item
  The colength of a simple A-module S is the dimension of the countermodule of S (over the field \(\operatorname{End}_{\mathrm{A}}(\mathrm{S})\) ). We say that an A-module M of finite length has finite colength if it admits a Jordan--Hölder series whose successive quotients are of finite colength; the sum of these colengths is called the colength of M and is denoted by \(\lambda(\mathrm{M})\) . If N is a submodule of M, then we have \(\lambda(\mathrm{M}) = \lambda(\mathrm{N}) + \lambda(\mathrm{M}/\mathrm{N})\) .
\end{enumerate}

Let A be an autoinjective ring (Exercise 8) such that the modules \(A_{s}\) and \(A_{d}\) are of finite colength. We say that A is a Frobenius ring if for every simple A-module S, we have \(\lambda(S) = \lambda(S^{*})\) .

\begin{enumerate}
\def\labelenumi{\alph{enumi})}
\item
  Let M be an A-module of finite length. If A is a Frobenius ring, then we have \(\lambda(M) = \lambda(M^{*})\) . If, moreover, A is an algebra of finite rank over a commutative field K, then we have \(\dim_K(M) = \dim_K(M^*)\) (use Exercise 7).
\item
  Let K be a commutative field, and let A be the subalgebra of \(\mathbf{M}_{6}(\mathrm{K})\) with basis the elements \(E_{11} + E_{55}\) , \(E_{12} + E_{56}\) , \(E_{21} + E_{65}\) , \(E_{22} + E_{66}\) , \(E_{33} + E_{44}\) , \(E_{13}\) , \(E_{23}\) , \(E_{45}\) , and \(E_{46}\) (where the \(E_{ij}\) are the matrix units of \(\mathbf{M}_{6}(\mathrm{K})\) ). Prove that A is an autoinjective algebra but not a Frobenius algebra (observe that the minimal left ideals of A are \(KE_{13} + KE_{23}\) , \(KE_{45}\) , and \(KE_{46}\) , and the minimal right ideals are \(KE_{13}\) , \(KE_{23}\) , and \(KE_{45} + KE_{46}\) ).
\end{enumerate}

\begin{enumerate}
\def\labelenumi{\arabic{enumi})}
\setcounter{enumi}{9}
\tightlist
\item
  Let \(\rho : \mathrm{A} \to \mathrm{B}\) be a ring homomorphism used to view B as a right A-module. Suppose that the left B-module \(\operatorname{Hom}_{\mathrm{A}}(\mathrm{B}, \mathrm{A})\) is isomorphic to \(\mathrm{B}_s\) .
\end{enumerate}

\begin{enumerate}
\def\labelenumi{\alph{enumi})}
\item
  Let M be a B-module, and let \(\varphi : \mathrm{B}_s \to \mathrm{Hom}_{\mathrm{A}}(\mathrm{B}, \mathrm{A})\) be an isomorphism. Prove that the mapping \(u \mapsto \varphi(1) \circ u\) from \(\mathrm{Hom}_{\mathrm{B}}(\mathrm{M}, \mathrm{B})\) to \(\mathrm{Hom}_{\mathrm{A}}(\mathrm{M}, \mathrm{A})\) is bijective.
\item
  Suppose, moreover, that the A-module B is finitely generated. The ring B is autoinjective (Exercise 8) if and only if the same holds for A.
\item
  Let G be a finite group. Show that the conditions above are satisfied when B is taken to be the A-module A{[}G{]} endowed with the ring structure such that \((ae_{g})(a^{\prime}e_{g^{\prime}})=aa^{\prime}e_{gg^{\prime}}\) for \(a,a^{\prime}\) in A and \(g,g^{\prime}\) in G. In particular, if K is a commutative field, then the algebra K{[}G{]} is autoinjective.
\end{enumerate}

